\documentclass[a4paper,12pt]{article}
\usepackage{amsmath}
\usepackage{graphicx}
\usepackage{multicol}
\usepackage{geometry}
\usepackage{fancyhdr}
\usepackage{tikz}
\usepackage{eso-pic}
\usepackage[spanish]{babel}

\geometry{top=1.5cm, bottom=1.5cm, left=1.5cm, right=1.5cm}

\pagestyle{fancy}
\fancyhf{}
\rhead{Esc. 4-117 Ejército de los Andes}
\lhead{Termodinámica y Máquinas Térmicas - Recuperatorio}
\cfoot{\thepage}

\begin{document}

\begin{center}
    \Large \textbf{Evaluación de Recuperatorio - Termodinámica y Máquinas Térmicas} \\
    \normalsize Duración: 80 minutos \\
    Alumno/a: \underline{\hspace{6cm}} \hspace{0.5cm} Fecha: \underline{\hspace{2cm}}
\end{center}

\vspace{0.3cm}

\begin{multicols}{2}

\section*{Problemas}

\begin{enumerate}

\item \textbf{Trabajo sobre gas a presión constante} \\
Un cilindro contiene $2\,\text{m}^3$ de gas que es comprimido a una presión constante de $5\,\text{atm}$. Calcule el trabajo necesario para realizar la compresión y explique brevemente la ecuación utilizada.

\vspace{0.3cm}

\item \textbf{Primer principio de la termodinámica} \\
Un gas realiza un trabajo de expansión de $3000\,\text{kgm}$, sin que cambie su energía interna. Determine la cantidad de calor absorbida o cedida por el sistema.

\vspace{0.3cm}

\item \textbf{Variación de energía interna} \\
Un pistón de área transversal $A = 0{,}012\,\text{m}^2$ se desplaza $d = 0{,}25\,\text{m}$ bajo una presión constante de $2{,}5\,\text{atm}$. Durante el proceso, el gas entrega $4{,}5\,\text{kcal}$ al ambiente. Calcule la variación de energía interna.

\vspace{0.3cm}

\item \textbf{Trabajo de circulación en compresor} \\
En una industria química, un compresor recibe $0{,}8\,\text{m}^3$ de aire a presión atmosférica y lo entrega comprimido a $0{,}03\,\text{m}^3$ y $3{,}5\,\text{atm}$, realizando un trabajo de compresión de $12000\,\text{J}$. Determine el trabajo de circulación total.

\vspace{0.3cm}

\item \textbf{Trabajo de circulación en contexto real} \\
Un sistema de aire acondicionado toma $1{,}2\,\text{m}^3$ de aire a presión atmosférica y lo comprime hasta $0{,}04\,\text{m}^3$ a una presión final de $2{,}8\,\text{atm}$. El trabajo de compresión realizado es de $14000\,\text{kgm}$. Calcule el trabajo de circulación total.

\end{enumerate}

\end{multicols}

\end{document}
